% Generated draft: requires mathematical and source review.
\documentclass[11pt]{article}
\usepackage[margin=25mm]{geometry}
\usepackage{fontspec}
\setmainfont{Latin Modern Roman}
\usepackage{amsmath,amssymb,mathtools,unicode-math}
\setmathfont{Latin Modern Math}
\usepackage{enumitem,needspace,xurl,hyperref}
\setlength{\parindent}{0pt}
\setlength{\parskip}{5pt}
\setlength{\emergencystretch}{3em}
\newcommand{\N}{\mathbb{N}}
\newcommand{\Z}{\mathbb{Z}}
\newcommand{\Q}{\mathbb{Q}}
\newcommand{\R}{\mathbb{R}}
\newcommand{\C}{\mathbb{C}}
\newcommand{\F}{\mathbb{F}}
\newcommand{\A}{\mathbb{A}}
\newcommand{\PP}{\mathbb{P}}
\newcommand{\E}{\mathbb{E}}
\newcommand{\eps}{\varepsilon}
\newcommand{\dd}{\,\mathrm{d}}
\begin{document}
\section*{UnsolvedMath ID: 10400109}
\textbf{Source problem code:} catalogue entry
\subsection{Definitions and mathematical statement}
The braid group $B_n$ has generators $\sigma_1,\ldots,\sigma_{n-1}$ with braid relations and commuting generators at distance at least two. In a common generic bracket normalization, let $A$ be indeterminate, $\delta=-A^2-A^{-2}$, and let $TL_n$ over $\Q(A)$ have generators $e_i$ with
\[
e_i^2=\delta e_i,\qquad e_ie_{i\pm1}e_i=e_i,\qquad e_ie_j=e_je_i\ (|i-j|\geq2).
\]
The braid representation takes values in its unit group:
\[
\rho_n(\sigma_i)=A\,1+A^{-1}e_i,\qquad
\rho_n(\sigma_i^{-1})=A^{-1}\,1+Ae_i.
\]
The question is whether $\rho_n$ is faithful, that is, $\ker\rho_n=\{1\}$, in the source's intended strand and parameter range. The displayed standard normalization must match that source before importing results. A generic indeterminate parameter and specialization at a number or root of unity are different problems; no specialization is silently chosen.

The supplied summary reports classical faithfulness for $n\leq3$, a simultaneous-faithfulness relation among the $B_4$ Temperley--Lieb/Jones/Burau representations from Bigelow, and a July 2026 preprint claiming faithfulness of $B_4$ Burau. Its stated conclusion of $B_4$ faithfulness is therefore recorded as dependent on that recent claim and the precise equivalence/normalization, not independently verified here. The summary leaves $n\geq5$ open and mentions a conjectural connection at five strands to Jones/HOMFLY unknot detection; no theorem-equivalence with that detection problem is asserted without its exact hypotheses.
\subsection{Short English statement}
Is the generic braid-group map into Temperley--Lieb units injective? Keep the strand range and parameter convention explicit, and treat the supplied recent four-strand solution claim with its actual source scope.
\subsection{Sources}
\begin{itemize}
\item Local UnsolvedMath registry, permanent ID 10400109, exact source code catalogue entry, Problem 6.1; source citation [188, Problem 3], identified in the summary as Kirby's Problem 3. No proposer/year in metadata. Registry status partially solved, classification PARTIAL-PROGRESS.
\item Ohtsuki, \emph{Problems on invariants of knots and 3-manifolds} (2002), Problem 6.1, PDF page 95, \url{https://msp.org/gtm/2002/04/gtm-2002-04-024s.pdf}. Supplied pdf-font-structure extraction accessed July 29, 2026; provenance distinguishes 231 bold numbered items. Bibliography entry [188] is absent; no PDF was accessed here.
\item Bigelow, ``Does the Jones polynomial detect the unknot?'', \emph{Journal of Knot Theory and Its Ramifications} \textbf{11} (2002), 493--505, \texttt{arXiv:math/0012086}, supplied for the $B_4$ simultaneous-faithfulness statement.
\item V. Bharathram, J. Birman and T. Brendle, \texttt{arXiv:2607.05283} (July 2026), supplied as claiming $B_4$ Burau faithfulness. No title or theorem text is given. Its correctness, status and applicability to the displayed normalization have not been checked. The summary's $n\leq4$ conclusion and five-strand detection connection are retained with these qualifications.
\item Provenance flags current primary-literature status and redistribution terms NEEDS\_REVIEW.
\end{itemize}
\end{document}
